\begin{abstract}
Machine-unlearning evaluations often compare an approximate intervention with a single
retrained model, despite stochastic training producing a population of counterfactual
models. This practice can obscure both natural training variability and a more basic
identifiability problem: a deletion request may not measurably change the learned function
under the selected assay. We present an evidence-traceable evaluation workflow that binds
deletion requests, retained sets, seeds, configurations, checkpoints, predictions, metrics,
and derived figures to hash-verified experiment records. The workflow represents exact
retraining with an ensemble and separates reference-assay sensitivity from subsequent
method evaluation.

We validate the workflow in P001, a retrospective pilot using a reduced Fashion-MNIST
split, one frozen one-percent IID deletion manifest, a small convolutional network, one
original training seed, and 15 independent exact-retraining seeds. The predefined
$k=10\rightarrow15$ seed-count diagnostic changed by \PilotJSConvergence{} for mean
Jensen--Shannon divergence and \PilotDisagreementConvergence{} for prediction
disagreement, below the frozen 10\% pilot threshold. Although three approximate methods
were executed, retrospective analysis found that the deletion itself was not separable
from exact-retraining variability under these two output metrics. Method proximity to the
exact ensemble therefore cannot be interpreted as evidence of successful removal in this
condition. High accuracy on IID-deleted examples is likewise not direct evidence of
memorization because held-out examples may remain classifiable from shared class structure.

P001 validates infrastructure integrity and initial reference-ensemble stability, not
certified unlearning, privacy equivalence, or generalization. A subsequent prospective
calibration found that JS divergence could identify full deletion of two tested classes,
but a candidate exact/exact-percentile equivalence margin failed to detect the frozen
one-percent residual-class violation. The evidence therefore supports distinguishability
analysis but not an equivalence claim. The resulting archive keeps the confirmatory method
matrix blocked until its decision rule is defensible.
\end{abstract}
